\section*{Bab \ref{ch:b1:plate-kinematics} --- Kinematika lempeng pada bola}

\begin{solution}{exo:b1:plate-kinematics:1}
(a) $\qty{0.210}{\degree\per\Ma} = 0.210\times\pi/180 = \qty{3.665e-3}{rad}$
per juta tahun, yaitu \qty{3.665e-9}{rad.yr^{-1}}; dengan
$\qty{1}{yr} = \qty{3.156e7}{s}$, \qty{1.161e-16}{rad.s^{-1}}. (b) $360/0.210 =
\qty{1714}{\Ma}$, lebih dari sepertiga umur Bumi. (c) $\omega R =
\num{3.665e-9}\times\qty{6.371e9}{mm} = \qty{23.35}{mm.yr^{-1}}$.
\end{solution}

\begin{solution}{exo:b1:plate-kinematics:2}
$v = \omega R\sin\Delta = 23.35\sin\Delta$: \qty{11.7}{mm.yr^{-1}} pada \ang{30},
\qty{20.2}{mm.yr^{-1}} pada \ang{60} dan \qty{23.35}{mm.yr^{-1}} pada \ang{90}. Penyebaran paling cepat di tempat batas paling jauh dari kutub, hingga
\ang{90}: antara Eurasia dan Amerika Utara, di ujung selatan batas,
dekat Azores, \ang{78} dari kutub.
\end{solution}

\begin{solution}{exo:b1:plate-kinematics:3}
Subkron berlangsung $1.072 - 0.988 = \qty{0.084}{\Ma}$, dan laju setengah dalam
\unit{mm.yr^{-1}} juga merupakan laju dalam \unit{km\per\Ma}. (a) $50\times0.084 =
\qty{4.2}{km}$ pada \qty{50}{mm.yr^{-1}}; $10\times0.084 = \qty{0.84}{km}$ pada
\qty{10}{mm.yr^{-1}}. (b) Pita itu mengapit sumbu:
$2\times10\times0.781 = \qty{15.6}{km}$.
\end{solution}

\begin{solution}{exo:b1:plate-kinematics:4}
Laju setengah adalah $66/2.581 = \qty{25.6}{mm.yr^{-1}}$ ke barat dan
$64/2.581 = \qty{24.8}{mm.yr^{-1}}$ ke timur; laju penuh adalah
$130/2.581 = \qty{50.4}{mm.yr^{-1}}$. Penyebaran
hampir simetris: sisi barat tumbuh sekitar
\qty{3}{\percent} lebih cepat.
\end{solution}

\begin{solution}{exo:b1:plate-kinematics:5}
$\cos\Delta = \sin\ang{61.8}\sin\ang{85} + \cos\ang{61.8}\cos\ang{85}
\cos(\ang{0} - \ang{139.6}) = 0.8813\times0.9962 + 0.4726\times0.08716\times
(-0.7615) = 0.87795 - 0.03137 = 0.84658$, sehingga $\Delta = \ang{32.16}$ dan
$v = 23.35\sin\ang{32.16} = \qty{12.4}{mm.yr^{-1}}$. Ini sedikit di atas
\qtyrange{11}{12}{mm.yr^{-1}} yang terukur, tetapi \omterm{def:b1:plate-kinematics:euler-pole}{vektor Euler} adalah satu kecocokan terbaik untuk 453 laju, dari mana laju-laju menyimpang rata-rata sebesar \qty{0.7}{mm.yr^{-1}}, dan laju-laju individual dekat Punggungan Gakkel yang lambat termasuk yang paling tidak presisi.
\end{solution}

\begin{solution}{exo:b1:plate-kinematics:6}
(a) $\cos\Delta = \sin\ang{61.8}\sin\ang{39.44} + \cos\ang{61.8}\cos\ang{39.44}
\cos(\ang{169.43}) = 0.20111$: $\Delta = \ang{78.40}$. (b) $v = 23.35
\sin\ang{78.40} = \qty{22.87}{mm.yr^{-1}}$, \qty{23}{mm.yr^{-1}} yang terukur
tepat di utara \omterm{def:b1:plate-kinematics:triple-junction}{titik tiga}. (c) Dengan rumus metode ini, pembilangnya
adalah $\sin\ang{169.43}\cos\ang{61.8} = 0.08668$ dan penyebutnya
$\cos\ang{39.44}\sin\ang{61.8} - \sin\ang{39.44}\cos\ang{61.8}\cos\ang{169.43}
= 0.97573$: kutub terlihat pada $\beta = \ang{5.08}$, hampir tepat utara, dan
Eurasia bergerak menuju $\ang{95.1}$. (d) Sesar transform sejajar dengan gerak
relatif: mereka membentang sekitar \ang{095}, hampir timur--barat.
\end{solution}

\begin{solution}{exo:b1:plate-kinematics:7}
(a) $\vec\omega = \omega(\cos\varphi_p\cos\lambda_p, \cos\varphi_p\sin\lambda_p,
\sin\varphi_p)$: $\vec\omega_{\mathrm{EU\,NA}} = (-0.0756, 0.0643, 0.1851)$ dan
$\vec\omega_{\mathrm{NB\,NA}} = (0.0333, 0.0282, 0.2289)$, dalam
\unit{\degree\per\Ma}. (b) $\vec\omega_{\mathrm{NB\,EU}} =
\vec\omega_{\mathrm{NB\,NA}} - \vec\omega_{\mathrm{EU\,NA}} = (0.1089, -0.0361,
0.0438)$, dengan panjang
\[
  \sqrt{0.011859 + 0.001303 + 0.001918} = \qty{0.1228}{\degree\per\Ma}.
\]
Kutubnya berada pada \omterm{def:b1:spherical-astronomy:geographic}{lintang}
$\arcsin(0.0438/0.1228) = \ang{20.9}$\,U dan \omterm{def:b1:spherical-astronomy:geographic}{bujur}
$\arctan(-0.0361/0.1089) = \ang{-18.3}$, yaitu \ang{18.3}\,B (komponen $x$
positif, tidak ada \ang{180} yang ditambahkan). (c) Batas antara Nubia dan Eurasia
memberikan sedikit pengukuran langsung: antara Azores dan Gibraltar, hanya \omterm{def:b1:spherical-astronomy:horizontal}{azimut} sesar
Gloria. Dua gerak penyebaran melintasi Atlantik diukur oleh ratusan laju, dan selisihnya
mewarisi presisinya.
\end{solution}

\begin{solution}{exo:b1:plate-kinematics:8}
(a) Sekitar $\vec e_z$ sebesar \ang{90}: $\cos\ang{90} = 0$, $\vec e_z\cdot\vec e_x = 0$
dan $\vec e_z\times\vec e_x = \vec e_y$, sehingga $\vec e_x$ pergi ke $\vec e_y$. Lalu
sekitar $\vec e_x$ sebesar \ang{90}: $\vec e_x\cdot\vec e_y = 0$ dan
$\vec e_x\times\vec e_y = \vec e_z$, sehingga $\vec e_y$ pergi ke $\vec e_z$: titik
berakhir di kutub utara. (b) Sekitar $\vec e_x$ dahulu: $\vec e_x$ berada pada sumbu,
$\vec e_x\times\vec e_x = \vec 0$ dan $(1 - \cos\ang{90})(\vec e_x\cdot\vec e_x)
\vec e_x = \vec e_x$, sehingga ia tidak bergerak; lalu sekitar $\vec e_z$ ia pergi ke
$\vec e_y$, pada \ang{0}\,U, \ang{90}\,T. (c) Hasilnya berbeda: \omterm{def:b1:plate-kinematics:finite-rotation}{rotasi berhingga}
tidak komutatif. Rekonstruksi berurutan harus digabung dengan mengalikan matriks rotasi dalam urutan yang benar, tidak pernah dengan menjumlahkan sudut
atau kutub.
\end{solution}

\begin{solution}{exo:b1:plate-kinematics:9}
(a) Dalam ruang kecepatan letakkan $a$ di titik asal, $b$ \qty{30}{mm.yr^{-1}} ke
timur dan $c$ \qty{40}{mm.yr^{-1}} ke selatan. (b) $\vec v_{CB} =
\vec v_{CA} - \vec v_{BA} = (0, -40) - (30, 0) = (-30, -40)$ (komponen timur, utara):
\qty{50}{mm.yr^{-1}} menuju \omterm{def:b1:spherical-astronomy:horizontal}{azimut}
$\arctan\bigl((-30)/(-40)\bigr) + \ang{180} = \ang{216.9}$, barat daya. (c) Punggungan
tegak lurus terhadap arah ini: ia membentang
$\ang{216.9} - \ang{90} = \ang{126.9}$, yaitu barat laut--tenggara.
\end{solution}

\begin{solution}{exo:b1:plate-kinematics:10}
(a) Dalam matriks Rodrigues, matriks identitas menyumbang $3\cos\alpha$, matriks
$\vec n\,\vec n^{\mathsf T}$ menyumbang jejak $n_x^2 + n_y^2 + n_z^2 = 1$, dikali
$1 - \cos\alpha$, dan matriks antisimetris tidak menyumbang apa-apa: $\operatorname{tr}R =
3\cos\alpha + 1 - \cos\alpha = 1 + 2\cos\alpha$. (b) Di sini $\operatorname{tr}R = 0$,
sehingga $\cos\alpha = -1/2$ dan $\alpha = \ang{120}$. Sumbunya memenuhi $R\vec n =
\vec n$: $R(x, y, z) = (z, x, y)$, sehingga $x = y = z$ dan $\vec n =
(1, 1, 1)/\sqrt3$, pada \omterm{def:b1:spherical-astronomy:geographic}{lintang} $\arcsin(1/\sqrt3) = \ang{35.26}$\,U dan
\omterm{def:b1:spherical-astronomy:geographic}{bujur} \ang{45}\,T. (c) $R\vec e_x = \vec e_y$, $R\vec e_y = \vec e_z$,
$R\vec e_z = \vec e_x$: ketiga sumbu saling bertukar, seperti yang dilakukan sepertiga putaran terhadap diagonal. Rumus Rodrigues dengan $\alpha = +\ang{120}$ terhadap $\vec n$
memang memberikan $\vec e_x\mapsto -\tfrac12\vec e_x + \tfrac12(1, 1, 1) +
\tfrac{\sqrt3}{2}\cdot\tfrac{1}{\sqrt3}(0, 1, -1) = (0, 1, 0)$, sehingga rotasinya
berlawanan arah jarum jam bila dilihat dari atas \ang{35.26}\,U, \ang{45}\,T.
\end{solution}

\begin{solution}{exo:b1:plate-kinematics:11}
(a) Pada sesar transform kecepatan relatif $\vec\omega\times\vec r$ berada sepanjang
sesar, sehingga sejajar dengan $\vec t$; ia tegak lurus terhadap $\vec\omega$, maka
$\vec t\cdot\vec\omega = 0$. Bidang yang tegak lurus $\vec t$ melalui pusat memuat $\vec r$ (karena $\vec t$ horizontal) dan $\vec\omega$: ia
adalah bidang \omterm{def:b1:spherical-astronomy:spherical-triangle}{lingkaran besar} yang melalui titik itu tegak lurus terhadap sesar,
dan kutubnya berada pada lingkaran itu. Dua bidang seperti itu bertemu sepanjang garis
$\vec t_1\times\vec t_2$. (b) $\vec t_1 = (0.5436, 0.8367, -0.0673)$ di \omterm{def:b1:plate-kinematics:triple-junction}{titik tiga}, dan
$\vec t_2 = (0.4942, 0.8639, -0.0968)$ di Islandia, sehingga
$\vec t_1\times\vec t_2 = (-0.02285, 0.01936, 0.05612)$, dengan panjang $0.06362$.
\omterm{def:b1:spherical-astronomy:geographic}{Lintangnya} adalah $\arcsin(0.05612/0.06362) = \ang{61.9}$\,U dan \omterm{def:b1:spherical-astronomy:geographic}{bujurnya}
$\ang{180} - \arctan(0.01936/0.02285) = \ang{139.7}$\,T (komponen $x$ negatif):
kutub dalam bab ini. (c) Buruk: kedua \omterm{def:b1:spherical-astronomy:spherical-triangle}{lingkaran besar} berpotongan pada sudut kecil, karena arah kedua transform hanya berbeda \ang{8} sepanjang \ang{26} batas.
Mengubah salah satu arah sebesar satu derajat menggeser kutub sekitar \ang{15}. Penentuan nyata menggabungkan banyak transform dengan laju penyebaran, yang variasinya sepanjang batas menetapkan jarak sudut ke kutub.
\end{solution}

\begin{solution}{exo:b1:plate-kinematics:12}
(a) $(94 - 4)/3 = \qty{30}{mm.yr^{-1}}$. (b) Selisih berjarak satu tahun adalah
$24 - 4 = 20$, $25 - 5 = 20$, $26 - 6 = 20$, $65 - 15 = 50$, $74 - 24 = 50$,
$75 - 25 = 50$, $76 - 26 = 50$, $85 - 65 = 20$ dan $94 - 74 = 20$ (mm). Diurutkan:
lima nilai 20 dan empat 50; median adalah \qty{20}{mm.yr^{-1}}. (c) Dua posisi
berjarak satu tahun berada pada fase yang sama dari osilasi musiman, yang saling meniadakan dalam selisihnya. Sebuah lompatan hanya memengaruhi pasangan yang melintasinya, di sini empat dari sembilan; selama kurang dari separuh pasangan melintasi lompatan, median mengabaikannya, sedangkan estimasi titik-ujung di (a) menganggap seluruh lompatan sebagai gerak dan menghasilkan nilai \qty{50}{\percent} terlalu tinggi. Metode ini gagal bila lompatan begitu sering sehingga sebagian besar pasangan satu tahun mengandung satu, atau bila stasiun bergerak tidak mantap, misalnya setelah gempa besar.
\end{solution}

\begin{solution}{pb:b1:plate-kinematics:1}
\textbf{1.} $\omega = 0.210\times(\pi/180)/\num{e6} = \qty{3.665e-9}{rad.yr^{-1}}$;
$\omega R = \num{3.665e-9}\times\qty{6.371e9}{mm} = \qty{23.35}{mm.yr^{-1}}$.

\textbf{2.} $\cos\Delta = \sin\ang{61.8}\sin\ang{64.5} + \cos\ang{61.8}
\cos\ang{64.5}\cos(\ang{-157.6}) = 0.79546 - 0.18809 = 0.60737$, sehingga
$\Delta = \ang{52.60}$.

\textbf{3.} $v = 23.35\sin\ang{52.60} = 23.35\times0.7944 =
\qty{18.55}{mm.yr^{-1}}$.

\textbf{4.} Pembilangnya adalah $\sin\ang{157.6}\cos\ang{61.8} = 0.18011$ dan
penyebutnya
\[
  \cos\ang{64.5}\sin\ang{61.8} - \sin\ang{64.5}\cos\ang{61.8}\cos\ang{157.6} =
  0.37940 + 0.39436 = 0.77376;
\]
keduanya positif, sehingga $\beta = \arctan 0.2328 =
\ang{13.1}$. Eurasia bergerak menuju $\ang{13.1} + \ang{90} = \ang{103.1}$.

\textbf{5.} Timur: $18.55\sin\ang{103.1} = \qty{18.07}{mm.yr^{-1}}$; utara:
$18.55\cos\ang{103.1} = \qty{-4.20}{mm.yr^{-1}}$, yaitu \qty{4.20}{mm.yr^{-1}}
ke selatan.

\textbf{6.} $\cos\Delta = 0.8851$, $\Delta = \ang{27.74}$, $v = 23.35
\sin\ang{27.74} = \qty{10.9}{mm.yr^{-1}}$. Kedua lempeng berputar terhadap kutub di
Siberia timur: Punggungan Gakkel dekat dengan kutub itu, dan kelajuannya turun dengan
$\sin\Delta$ meskipun lempengnya sama.

\textbf{7.} Barat: $94/9.779 = \qty{9.61}{mm.yr^{-1}}$; timur: $97/9.779 =
\qty{9.92}{mm.yr^{-1}}$.

\textbf{8.} $(94 + 97)/9.779 = 191/9.779 = \qty{19.53}{mm.yr^{-1}}$.

\textbf{9.} Dari C2An: $54.1/2.581 = \qty{20.96}{mm.yr^{-1}}$; dari dasar
C1n: $19.8/0.781 = \qty{25.35}{mm.yr^{-1}}$. Keduanya tidak sepakat: semakin muda batasnya, semakin tinggi lajunya. Perubahan laju yang nyata mungkin terjadi, tetapi selisih jarak yang tetap, sama di semua umur, akan menghasilkan persis hal ini.

\textbf{10.} Rasio jarak barat terhadap timur adalah 0.96 hingga 0.98: sisi timur tumbuh beberapa persen lebih banyak. Laju penuh hanya bergantung pada
jumlah kedua jarak, sehingga asimetri tidak mengubahnya.

\textbf{11.} $(191 - 5.0)/9.779 = \qty{19.02}{mm.yr^{-1}}$;
$(54.1 - 5.0)/2.581 = \qty{19.02}{mm.yr^{-1}}$; $(19.8 - 5.0)/0.781 =
\qty{18.95}{mm.yr^{-1}}$. Ketiganya kini sepakat dalam
\qty{0.1}{mm.yr^{-1}}. Pergeseran \qty{5}{km} yang sama adalah seperempat dari \qty{19.8}{km}
pasangan termuda tetapi kurang dari \qty{3}{\percent} dari \qty{191}{km} pasangan
tertua; batas-batas tua memberikan laju yang lebih andal.

\textbf{12.} Laju setengah yang terkoreksi adalah $186/9.779/2 = \qty{9.51}{mm.yr^{-1}}$,
dan $9.51\times0.084 = \qty{0.80}{km}$, sepertiga dari \qty{2.5}{km} tempat
setiap batas tergeser. Orang akan mengharapkan pita yang sempit itu kabur
karenanya. Bahwa ia dikenali di Punggungan Reykjanes, sedangkan tidak ada di
Punggungan Gakkel, yang kekurangan magma, menunjukkan bahwa pasokan magma yang melimpah
mencatat medan dengan setia bahkan dalam pita yang sempit.

\textbf{13.} Laju terkoreksi, 18.95 hingga \qty{19.02}{mm.yr^{-1}}, sepakat dengan
\qty{19.0}{mm.yr^{-1}} vektor lempeng; tanpa koreksi, mereka melebihinya sebesar 3 hingga 33 persen.

\textbf{14.} $\cos\Delta = \sin\ang{64.139}\sin\ang{59.018} +
\cos\ang{64.139}\cos\ang{59.018}\cos\ang{27.554} = 0.97054$, sehingga $\Delta =
\ang{13.94}$; panjang garis itu $2R\sin(\Delta/2) = 12742\times\sin\ang{6.97} =
\qty{1546}{km}$.

\textbf{15.} $1000/17.5 = 57$ tahun.

\textbf{16.} Kedua arah berbeda $\ang{123.1} - \ang{114.8} =
\ang{8.3}$: komponen horizontal sepanjang \omterm{def:b1:spherical-astronomy:spherical-triangle}{lingkaran besar} adalah
$18.9\cos\ang{8.3} = \qty{18.70}{mm.yr^{-1}}$, dan sepanjang garis
$18.70\cos\ang{6.97} = \qty{18.56}{mm.yr^{-1}}$, sekitar \qty{18.6}{mm.yr^{-1}}.

\textbf{17.} Garis itu bertambah \qty{17.5}{mm.yr^{-1}}, \qty{6}{\percent}
lebih sedikit daripada \qty{18.6}{mm.yr^{-1}} yang diprediksi. Kedua stasiun berada jauh dari batas lempeng, sehingga selisihnya bukan lokal. Entah lempeng bergerak sedikit lebih lambat hari ini daripada rata-rata selama tiga juta tahun terakhir, atau salah satu dari dua estimasi membawa galat sistematis kecil; pengukuran saja tidak dapat memberitahukannya. \omterm{def:b1:plate-kinematics:euler-pole}{Vektor Euler} yang dihitung dari penentuan posisi satelit saja memprediksi \qty{17.1}{mm.yr^{-1}} untuk garis ini, dekat dengan nilai terukur.

\textbf{18.} Timur: $21.12 - (-1.07) = \qty{22.19}{mm.yr^{-1}}$; utara:
$-6.97 - (-0.97) = \qty{-6.00}{mm.yr^{-1}}$. Kelajuan $\sqrt{22.19^2 + 6.00^2} =
\qty{23.0}{mm.yr^{-1}}$, menuju $\ang{180} - \arctan(22.19/6.00) =
\ang{105.1}$.

\textbf{19.} Arahnya dekat dengan \ang{103} dari vektor lempeng, tetapi kelajuannya melebihi \qty{18.6}{mm.yr^{-1}} sebesar seperempat. Kedua stasiun berada di dalam zona tempat batas berdeformasi, dan Höfn, dekat tudung es terbesar, bergerak mengikuti pemantulan \omterm{def:b1:layered-earth:layers}{kerak} ketika es menipis: naik sebesar \qty{12.9}{mm.yr^{-1}}, dan ke samping, menjauhi es. Stasiun di dalam zona batas mengukur gerak lempeng ditambah deformasi lokal.

\textbf{20.} $0.210\times10 = \ang{2.10}$.

\textbf{21.} Titik itu menempuh \omterm{def:b1:spherical-astronomy:spherical-triangle}{lingkaran kecil} berjari-jari $R\sin\Delta$:
$6371\times0.7944\times(2.10\times\pi/180) = \qty{185.5}{km}$. Lebar \omterm{def:b1:layered-earth:layers}{kerak} yang terkoreksi
yang terbentuk sejak tepi muda C5n, \qty{186}{km} dalam
\qty{9.8}{Ma}, hampir sama: gerak itu telah mantap selama sepuluh
juta tahun.

\textbf{22.} Komponen (timur, utara): Eurasia
$(22.9\sin\ang{95.1}, 22.9\cos\ang{95.1}) = (22.81, -2.04)$, Nubia
$(19.2\sin\ang{103.8}, 19.2\cos\ang{103.8}) = (18.65, -4.58)$. Eurasia relatif
terhadap Nubia: $(4.16, 2.54)$, yaitu \qty{4.9}{mm.yr^{-1}} menuju
$\arctan(4.16/2.54) = \ang{58.6}$, timur-timur laut.

\textbf{23.} Sebuah \omterm{def:b1:plate-kinematics:triple-junction}{titik tiga} dari tiga punggungan \omterm{def:b1:plate-kinematics:triple-junction}{stabil} apa pun kecepatannya. Kecepatan pada soal 22 membuat Eurasia dan Nubia bergerak saling menjauh, dari timur-timur laut ke barat-barat daya: batas mereka di titik itu harus membuka, perlahan; itu adalah Rift Terceira, yang arah timur-tenggaranya membuat pembukaan itu miring.

\textbf{24.} Bagian I: \qty{18.6}{mm.yr^{-1}} di pusat Islandia (\qty{19.0}{mm.yr^{-1}} di profil Bagian II). Bagian II: 18.95 hingga \qty{19.02}{mm.yr^{-1}} setelah dikoreksi. Bagian III: garis Reykjavik--Stavanger bertambah \qty{17.5}{mm.yr^{-1}}, di mana vektor lempeng memprediksi \qty{18.6}{mm.yr^{-1}}; pasangan di dalam Islandia memberikan \qty{23.0}{mm.yr^{-1}}, yang bukan gerak lempeng. Vektor lempeng dan pita sepakat; satelit berbeda sebesar \qty{6}{\percent}. Pita merata-ratakan selama hingga sepuluh juta tahun, vektor lempeng selama \qty{3.16}{Ma}; satelit mengukur gerak dua puluh hingga tiga puluh tahun terakhir.

\textbf{25.} Vektor lempeng memberikan sekitar \qty{18.55}{mm.yr^{-1}} di seluruh Islandia, dan pita terkoreksi \qty{19.0}{mm.yr^{-1}} di sebelahnya: hingga milimeter per tahun terdekat, \omterm{def:b1:plate-kinematics:spreading-rate}{laju penyebaran penuh} di seluruh Islandia adalah
\textbf{\qty[text-series-to-math=true]{19}{mm.yr^{-1}}}.
\end{solution}
